% !TEX root = template-iken-lua.tex
% !TEX recipe = LuaLaTeX (kian-lua)
% ============================================================
%  意見書テンプレート  (kian-v08 LuaLaTeX 版)
%  コンパイル: latexmk -r latexmkrc-lua template-iken-lua.tex
%              または VS Code の「LuaLaTeX (kian-lua)」レシピ
%  プレースホルダー: 【...】を実際の内容に置き換えてください
%  異体字: \ivs{文字} で IPAmj 明朝字形に切り替え（例: \ivs{角󠄁}）
%  ※ \SecA{} 等の {} 内に % を書くとコメントになりエラーが出ます
%  見出し番号: \SecZ=番号なし \SecA=第１ \SecB=１ \SecC=⑴ \SecD=ア \SecE=(ア)
%  ※ 見出し（\SecA{} 等）は本文なしで連続させると番号表示が崩れるため、
%    各見出しの直後に【ここに記載】のプレースホルダー本文を入れています
% ============================================================

\documentclass{ltjsarticle}
\usepackage{kian-lua}                % kian-lua/ を TEXINPUTS に含めてコンパイル
% \usepackage{kian-ivs,luatexja-ruby}  % IVS 異体字（\ivs{文字}）・ルビ用

\Idcs{
  [o]{R８(　)　　:【事件名】\\原告:【氏/名】:\\被告:【氏/名】}
}

\Tit{意見書}
\Dat{令和８年　　月　　日}

%\senderhlength{18}{4}{2}{10}% \senderhlength{①}{②}{③}{④}
% ①左端からのインデント（開始位置）② （肩書幅）「弁護士」という肩書を配置するボックスの横幅です。③ （肩書と名前の間隔）「弁護士」と「名前」の間に空けるスペースの幅です。④（名前幅）「【氏名】」という名前を均等割り付け（配置）するボックスの横幅です。
%  意見書作成者（専門家）を \Sens で複数行表示
%  氏名に異体字がある場合: \ivs{角󠄁}田　一郎
\Sens{
  [0.5]{%
    【所属機関・肩書】/【氏　名】\\
  }
}

% ============================================================

\begin{document}

\maketitle

% ---------------------------------------------------------- %
%  意見の依頼
% ---------------------------------------------------------- %

\SecA{意見の依頼}



% ---------------------------------------------------------- %
%  意見書作成者の略歴（専門家意見書の場合）
% ---------------------------------------------------------- %

\SecA{意見書作成者の略歴}



% ---------------------------------------------------------- %
%  事案の概要
% ---------------------------------------------------------- %

\SecA{事案の概要}



% ---------------------------------------------------------- %
%  前提事実
% ---------------------------------------------------------- %

\SecA{前提事実}

\SecB{}

【ここに記載】

\SecB{}

【ここに記載】

% ---------------------------------------------------------- %
%  意見
% ---------------------------------------------------------- %

\SecA{意見}

\SecB{}

【ここに記載】

\SecB{}

【ここに記載】

% ---------------------------------------------------------- %
%  結論
% ---------------------------------------------------------- %

\SecA{結論}



\vspace{2\zw}

\begin{flushright}
以　上
\end{flushright}

\end{document}
